% !PRINTED 596
\weberpage{596}
Setzt man $h = p$, so ergiebt sich hieraus
\[ \tau\mkern1mu|\mkern1mu\chi_1^p \equiv \tau\ (\mathrm{mod}\ \mathfrak{p}^3), \]
und auf die gleiche Weise schliesst man hieraus für jeden Exponenten $\nu$
\weq{17}{\tau\mkern1mu|\mkern1mu\chi_1^{p^\nu} \equiv \tau\ (\mathrm{mod}\ \mathfrak{p}^{\nu+2}).}
Um diese Formel allgemein zu beweisen, wendet man die vollständige Induction an. Man nimmt (17) als bewiesen an und setzt
\[ \tau\mkern1mu|\mkern1mu\chi_1^{p^\nu} = \tau + \pi^{\nu+2}\alpha; \]
daraus durch wiederholte Anwendung
\[
\begin{aligned}
\tau\mkern1mu|\mkern1mu\chi_1^{2p^\nu} &\equiv \tau + 2\pi^{\nu+2}\alpha \\
\tau\mkern1mu|\mkern1mu\chi_1^{3p^\nu} &\equiv \tau + 3\pi^{\nu+2}\alpha\ (\mathrm{mod}\ \mathfrak{p}^{\nu+3}) \\
&\cdots\cdots\cdots\cdots \\
\tau\mkern1mu|\mkern1mu\chi_1^{p^{\nu+1}} &\equiv \tau + p\pi^{\nu+2}\alpha,
\end{aligned}
\]
also die Formel (17) für $\nu + 1$.

Nun kann man $\nu$ immer so gross annehmen, dass von den Differenzen
\[ \tau - \tau_1,\ \tau - \tau_2,\ \tau - \tau_3,\ \ldots, \]
wenn $\tau$ von $\tau_1, \tau_2, \tau_3, \ldots$ verschieden ist, keine durch $\mathfrak{p}^{\nu-2}$ theilbar ist, und dann folgt aus (17), dass $\chi_1^{p^\nu}$ die identische Substitution sein muss. \gesperrt{Es ist hiernach der Grad eines jeden Elementes $\chi_1$ eine Potenz von $p$, und folglich kann nach dem Cauchy-Sylow'schen Satze (\S.\,29, I.) im Grade von $X^{(1)}$ keine von $p$ verschiedene Primzahl aufgehen, wie bewiesen werden sollte.}

Wenn also $g$ nicht durch $p$ theilbar ist, dann ist sicher $X^{(1)}$ die Einheitsgruppe.

Man kann auf diese Weise fortfahren und eine Gruppe $X^{(2)}$ vom Grade $g_2$ bilden aus der Congruenz
\[ \tau\mkern1mu|\mkern1mu\chi_2 \equiv \tau\ (\mathrm{mod}\ \mathfrak{p}^3), \]
$X^{(2)}$ ist ein Theiler von $X^{(1)}$, also $g_2$ gleichfalls eine Potenz von $p$, und so muss man schliesslich zur Einheitsgruppe gelangen. Unter den auf einander folgenden Gruppen $X^{(1)}, X^{(2)}, \ldots$ können unter Umständen auch mehrere einander gleich sein\footnote{\gesperrt{Hilbert} nennt $\Psi$ die Zerlegungsgruppe, $X$ die Trägheitsgruppe, $X^{(1)}$ die Verzweigungsgruppe des Ideals $\mathfrak{p}$. In der oben erwähnten Abhandlung ist noch bewiesen, dass $g_1$ die höchste in $g$ aufgehende Potenz von $p$ ist, dass $X^{(1)}$ ein Normaltheiler von $X$ ist, und dass die ganze Gruppe $\Psi$ metacyklisch ist.}.
